% Part: first-order-logic
% Chapter: tableaux
% Section: proof-theoretic-notions

\documentclass[../../../include/open-logic-section]{subfiles}

\begin{document}

\iftag{FOL}
      {\olfileid{fol}{tab}{ptn}}
      {\olfileid{pl}{tab}{ptn}}
      
\olsection{सिद्धता-उपपत्तीय संकल्पना}

\begin{editorial}
या विभागात टॅब्लोसाठी सिद्धतायोग्यता संबंध आणि सुसंगतता यांच्या व्याख्या
एकत्रित केल्या आहेत.
\end{editorial}

\begin{explain}
जशा आपण अनेक महत्त्वाच्या चिन्हार्थविषयक संकल्पना (वैधता, तार्किक निष्पन्नता,
पूर्ततायोग्यता) परिभाषित केल्या आहेत, तशाच त्यांना अनुरूप
\emph{सिद्धता-उपपत्तीय संकल्पना} आता परिभाषित करू. !!{structure}s मध्ये
!!{sentence}s ची पूर्ति होते की नाही याचा आधार घेऊन या संकल्पना परिभाषित
केल्या जात नाहीत; त्याऐवजी ठरावीक बंद !!{tableau}x च्या अस्तित्वाचा आधार
घेतला जातो. या दोन प्रकारच्या संकल्पना एकमेकांशी जुळतात, हा एक महत्त्वाचा
शोध होता. त्या जुळतात, हाच \emph{निर्दोषता प्रमेय} आणि
\emph{संपूर्णता प्रमेय} यांचा आशय आहे.
\end{explain}


\begin{defn}[प्रमेये]
एखादे !!{sentence}~$!A$ हे \emph{प्रमेय} असते, जर
$\sFmla{\False}{!A}$ साठी बंद !!{tableau} असेल. $!A$ हे प्रमेय असेल
तर आपण $\Proves !A$ लिहितो आणि ते प्रमेय नसेल तर $\Proves/ !A$ लिहितो.
\end{defn}

\begin{defn}[!!^{derivability}]
!!^a{sentence} $!A$ हे !!{sentence}s च्या संच~$\Gamma$ \emph{पासून
!!{derivable}} असते, म्हणजे $\Gamma \Proves !A$, तेव्हा आणि केवळ तेव्हाच,
जेव्हा एखादा सांत संच $\{!B_1, \dots, !B_n\} \subseteq \Gamma$ असा असतो
आणि पुढील संचासाठी बंद !!{tableau} असतो:
\[
\{
  \sFmla{\False}{!A},
  \sFmla{\True}{!B_1}, \dots,
  \sFmla{\True}{!B_n}
  \}.
\]
$!A$ हे $\Gamma$ पासून !!{derivable} नसेल, तर आपण $\Gamma \Proves/
!A$ लिहितो.
\end{defn}

\begin{defn}[सुसंगतता]
!!{sentence}s चा संच~$\Gamma$ हा \emph{विसंगत} असतो तेव्हा आणि केवळ तेव्हाच,
जेव्हा एखादा सांत संच $\{!B_1, \dots, !B_n\} \subseteq \Gamma$ असा असतो
आणि पुढील संचासाठी बंद !!{tableau} असतो:
\[
\{
  \sFmla{\True}{!B_1}, \dots,
  \sFmla{\True}{!B_n}  
  \}.
\]
$\Gamma$ विसंगत नसेल, तर त्याला \emph{सुसंगत} म्हणतो.
\end{defn}

\begin{prop}[परावर्तिता]
\ollabel{prop:reflexivity}
$!A \in \Gamma$ असेल, तर $\Gamma \Proves !A$.
\end{prop}

\begin{proof}
$!A \in \Gamma$ असेल, तर $\{!A\}$ हा $\Gamma$ चा सांत उपसंच आहे आणि पुढील !!{tableau}
\begin{oltableau}
  [\sFmla{\False}{\formula{A}}, just = \TAss
    [\sFmla{\True}{\formula{A}}, just = \TAss,close]
  ]
\end{oltableau}
बंद आहे.
\end{proof}
  

\begin{prop}[एकस्वनिकता]
\ollabel{prop:monotonicity}
$\Gamma \subseteq \Delta$ आणि $\Gamma \Proves !A$ असेल, तर $\Delta
\Proves !A$.
\end{prop}

\begin{proof}
$\Gamma$ चा कोणताही सांत उपसंच हा $\Delta$ चासुद्धा सांत उपसंच असतो.
\end{proof}

\begin{prop}[संक्रमकता]
\ollabel{prop:transitivity}
$\Gamma \Proves !A$ आणि $\{!A\} \cup
\Delta \Proves !B$ असेल, तर $\Gamma \cup \Delta \Proves !B$.
\end{prop}

\begin{proof}
$\{!A\} \cup \Delta \Proves !B$ असेल, तर $\Delta_0 =
\{!C_1, \dots, !C_n\} \subseteq \Delta$ असा सांत उपसंच असतो आणि पुढील संचासाठी
\begin{align*}
\{\sFmla{\False}{!B}, & \sFmla{\True}{!A}, \sFmla{\True}{!C_1},
\dots, \sFmla{\True}{!C_n}\}
\intertext{बंद !!{tableau} असतो. जर $\Gamma \Proves !A$ असेल, तर
  $\{!D_1, \dots, !D_m\} \subseteq \Gamma$ असा संच असतो आणि पुढील संचासाठीही}
\{\sFmla{\False}{!A}, & \sFmla{\True}{!D_1},
\dots, \sFmla{\True}{!D_m}\}
\end{align*}
बंद !!{tableau} असतो.

आता पुढील गृहीतके असलेला !!{tableau} विचारात घ्या:
\[
\sFmla{\False}{!B},
\sFmla{\True}{!C_1}, \dots, \sFmla{\True}{!C_n},
\sFmla{\True}{!D_1}, \dots, \sFmla{\True}{!D_m}.
\]
त्यात~$!A$ वर \Cut{} नियम लावा. यामुळे दोन शाखा तयार होतात: एका शाखेत
$\sFmla{\True}{!A}$, तर दुसऱ्या शाखेत $\sFmla{\False}{!A}$ असते. त्यामुळे
पहिल्या शाखेत पुढील सर्व गृहीतके उपलब्ध असतात:
\[
\{\sFmla{\False}{!B}, \sFmla{\True}{!A},
\sFmla{\True}{!C_1}, \dots, \sFmla{\True}{!C_n}\}
\]
या गृहीतकांसाठी बंद !!{tableau} असल्यामुळे तो त्या शाखेला जोडता येतो;
$\sFmla{\True}{!A}$ मधून जाणारी प्रत्येक शाखा बंद होते. दुसऱ्या शाखेत
पुढील सर्व गृहीतके उपलब्ध असतात:
\[
\{\sFmla{\False}{!A}, \sFmla{\True}{!D_1}, \dots,
\sFmla{\True}{!D_m}\}
\]
म्हणून दुसरी शाखासुद्धा पूर्ण करून बंद !!{tableau} मिळवता येतो. यावरून
$\Gamma \cup \Delta \Proves !B$ हे दिसते.
\end{proof}

विशेषतः, याचा अर्थ असा की $\Gamma \Proves !A$ आणि $!A
\Proves !B$ असेल, तर $\Gamma \Proves !B$. तसेच $!A_1,
\dots, !A_n \Proves !B$ असेल आणि प्रत्येक~$i$ साठी $\Gamma \Proves !A_i$
असेल, तर $\Gamma \Proves !B$, हेसुद्धा यावरून मिळते.

\begin{prop}
\ollabel{prop:incons}
$\Gamma$ विसंगत असतो तेव्हा आणि केवळ तेव्हाच, जेव्हा प्रत्येक
  !!{sentence}~$!A$ साठी $\Gamma \Proves !A$.
\end{prop}

\begin{proof}
सराव.
\end{proof}

\tagprob{FOL}
\begin{prob}
\olref[fol][tab][ptn]{prop:incons} सिद्ध करा.
\end{prob}
\tagendprob

\tagprob{notFOL}
\begin{prob}
\olref[pl][tab][ptn]{prop:incons} सिद्ध करा.
\end{prob}
\tagendprob

\begin{prop}[संहतता]
  \ollabel{prop:proves-compact}
  \begin{enumerate}
  \item $\Gamma \Proves !A$ असेल, तर एखादा सांत उपसंच $\Gamma_0
    \subseteq \Gamma$ असा आहे की $\Gamma_0 \Proves !A$.
  \item $\Gamma$ चा प्रत्येक सांत उपसंच सुसंगत असेल, तर $\Gamma$ सुसंगत
    असतो.
  \end{enumerate}
\end{prop}

\begin{proof}
  \begin{enumerate}
    \item $\Gamma \Proves !A$ असेल, तर $\Gamma_0 =
      \{!B_1, \dots, !B_n\}$ असा त्याचा सांत उपसंच आणि पुढील संचासाठी बंद !!{tableau} असतो:
      \[
      \{\sFmla{\False}{!A}, \sFmla{\True}{!B_1}, \dots, \sFmla{\True}{!B_n}\}
      \]
      हा !!{tableau} $\Gamma_0 \Proves !A$ हेसुद्धा दाखवतो.
    \item $\Gamma$ विसंगत असेल, तर
      $\Gamma_0 = \{!B_1, \dots, !B_n\}$ असा त्याचा एखादा सांत उपसंच असून पुढील
      संचाचा बंद !!{tableau} असतो:
      \[
      \{\sFmla{\True}{!B_1}, \dots, \sFmla{\True}{!B_n}\}
      \]
      हा बंद !!{tableau} $\Gamma_0$ विसंगत असल्याचे दाखवतो.
  \end{enumerate}
\end{proof}

\end{document}
